\section{Finite relational witnesses and shrinking viability}
\label{sec:relational}

Non-unique outputs are not an edge case.  They arise whenever roots collide,
eigenvectors rotate inside an eigenspace, factorizations reorder, local charts
overlap, minimizers are not isolated, or a finite proof has several equally
valid certificates.  Requiring a coherent exact selector can add logical
strength unrelated to the theorem.  The relational layer avoids that error.

\subsection{Finite cloud towers}

For levels \(i=0,\ldots,H\), let \(C_i\) be a finite coded cloud.  Its
candidates are bounded by a PR term \(B_i\), validity is a decidable predicate
\(V_i(a)\), and compatibility between adjacent levels is a decidable relation
\(E_i(a,b)\).  Think of level as precision, truncation order, modulus, or stage
of a finite cover.

Define backward viability relative to horizon \(H\):
\begin{align*}
 W_H^H(a)&\defeq V_H(a),\\
 W_i^H(a)&\defeq V_i(a)\wedge
   \exists b<B_{i+1}\,
   \bigl(V_{i+1}(b)\wedge E_i(a,b)\wedge W_{i+1}^H(b)\bigr).
\end{align*}
All existential quantifiers are bounded.  The predicate is therefore compiled
by a reverse-time primitive recursion using
Theorem~\ref{thm:bounded-booleanization}.

\begin{theorem}[Viability is finite path existence]
\label{thm:viability-path}
For every \(i\leq H\) and candidate \(a\in C_i\), \(W_i^H(a)\) holds iff
there is a sequence
\[
 a_i=a,a_{i+1},\ldots,a_H
\]
with \(V_j(a_j)\) for all \(j\) and
\(E_j(a_j,a_{j+1})\) for all \(j<H\).
\end{theorem}

\begin{proof}
Use induction on the remaining horizon \(H-i\).  At zero, the claimed path is
the singleton \((a)\), so the statement is exactly \(V_H(a)\).  At a
successor, the defining clause for \(W_i^H(a)\) supplies a compatible valid
\(b\) with \(W_{i+1}^H(b)\); the induction hypothesis extends \(b\) to a
path.  Conversely, take the second element of any path and apply the induction
hypothesis to its tail.  The proof is constructive and finite.
\end{proof}

In iPRA, the induction invariant is a Boolean equation comparing the viability
checker with a bounded path-checker.  Both are PR.  No formula quantifies over
an actual function or infinite sequence.

\subsection{Least finite paths}

Assume some initial candidate is viable.  Apply least bounded search to choose
the least viable \(a_0\).  Given \(a_i\), choose the least \(a_{i+1}\) that is
compatible and viable.  The viability clause guarantees existence at every
step.

\begin{theorem}[Finite dependent choice by bounded search]
\label{thm:finite-choice}
If \(\exists a<B_0\,W_0^H(a)\), there is a primitive-recursive function that
returns the lexicographically least valid compatible path through horizon
\(H\).  iPRA proves every row's validity, every edge's compatibility, and the
lexicographic minimality of the path.
\end{theorem}

\begin{proof}
The path constructor is a bounded primitive recursion on the level.  At each
step use Theorem~\ref{thm:least-search} on the finite successor cloud, with the
predicate \(E_i(a_i,b)\wedge W_{i+1}^H(b)\).  The viability theorem supplies
existence.  The least-search theorem supplies validity, compatibility, and
minimality.  A table fold accumulates the path code.
\end{proof}

This theorem is sometimes described informally as dependent choice, but it is
not an infinitary choice principle.  The horizon and every candidate range are
finite and explicitly bounded.

\subsection{Why an infinite branch is unnecessary}

Many applications need an approximation at each requested precision, not a
single infinite branch object.  For precision \(p\), choose a PR horizon
\(K(p)\geq p\), compute the least path through that finite horizon, and return
its level-\(p\) observation.  A separate quantitative theorem proves coherence
as \(p\) changes.

This approach is stronger computationally than postulating an infinite branch:
it gives a primitive-recursive procedure for every finite observation and
keeps the proof inside the bounded fragment.

\subsection{Shrinking viability theorem}

Let \((M,d)\) be a coded rational metric presentation.  Each candidate
\(a\in C_i\) has an observation \(O_i(a)\in M\).  For horizon \(H\geq i\), let
\[
 S_i^H=\{a\in C_i:W_i^H(a)\}.
\]
Assume the following quantitative data.

\begin{enumerate}
\item \emph{Nonemptiness}: for every required horizon, \(S_0^H\neq\varnothing\).
\item \emph{Shrinking}: for all \(H\geq i\),
\[
 \diam O_i(S_i^H)\leq\delta_i.
\]
\item \emph{Edge motion}: if \(E_i(a,b)\), then
\[
 d(O_i(a),O_{i+1}(b))\leq\eta_i.
\]
\item \emph{Tail modulus}: a PR function \(M(p)\) satisfies
\[
 \sum_{i=M(p)}^{m-1}(\delta_i+\eta_i)\leq 2^{-p}
 \quad\text{for all }m>M(p).
\]
The finite sums are rational and their bound is proved by a bounded checker or
closed estimate.
\end{enumerate}

For a nondecreasing horizon schedule \(K(i)\geq i+1\), let \(P^H\) be the
least path through horizon \(H\), and define
\[
 a_i=P^{K(i)}_i,\qquad x_i=O_i(a_i).
\]

\begin{theorem}[Shrinking finite-horizon coherence]
\label{thm:shrinking-viability}
Under the preceding hypotheses,
\[
 d(x_i,x_{i+1})\leq\delta_i+\eta_i.
\]
Consequently, for \(m>n\),
\[
 d(x_n,x_m)\leq\sum_{i=n}^{m-1}(\delta_i+\eta_i),
\]
and \((x_i)\) is a primitive-recursive Cauchy name with modulus \(M\).
\end{theorem}

\begin{proof}
Assume \(K(i+1)\geq K(i)\); replace the schedule by its running maximum if
necessary.  Let \(c_i=P^{K(i+1)}_i\).  Both \(a_i\) and \(c_i\) extend to
horizon \(K(i)\): the first by definition and the second because its longer
path can be truncated.  Hence they lie in \(S_i^{K(i)}\), and shrinking gives
\[
 d(O_i(a_i),O_i(c_i))\leq\delta_i.
\]
The path \(P^{K(i+1)}\) contains the edge
\(c_i\to a_{i+1}\), so edge motion gives
\[
 d(O_i(c_i),O_{i+1}(a_{i+1}))\leq\eta_i.
\]
The triangle inequality proves the one-step bound.  Telescoping proves the
finite-distance bound, and the supplied tail modulus proves the Cauchy
condition.
\end{proof}

The theorem uses no compactness principle and no infinite path.  At each
precision it computes one finite least path.  Coherence follows from shrinking
of the entire viable cloud, not from persistence of an exact witness.

\subsection{Nested viable sets and a simpler special case}

In some presentations all level-\(i\) candidates valid to a sufficiently long
horizon already have diameter \(\delta_i\), and adjacent observations are
identical under restriction, so \(\eta_i=0\).  Then
Theorem~\ref{thm:shrinking-viability} reduces to the familiar statement that
nested finite approximation clouds with shrinking diameter determine a
regular name.  The finite-horizon proof is preferable to an appeal to the
Cantor intersection theorem because it remains first-order and quantitative.

\subsection{Exact labels from a positive gap}

Let \(L=\{\lambda_0,\ldots,\lambda_{r-1}\}\) be a finite list of exact coded
labels with pairwise separation
\[
 d(\lambda_i,\lambda_j)\geq g>0\qquad(i\neq j).
\]
Suppose an interval or ball \(B\) has diameter less than \(g\) and is proved to
contain some label.

\begin{lemma}[Gap collapse]
\label{lem:gap-collapse}
The ball \(B\) contains exactly one label.  Bounded least search returns its
index, and iPRA proves that index is stable under every refinement
\(B'\subseteq B\) that is also proved to contain a label.
\end{lemma}

\begin{proof}
Existence is assumed.  If two distinct labels lay in \(B\), their distance
would be at most the diameter of \(B\), contradicting the separation bound.
The unique index is therefore the least index satisfying the bounded
containment predicate.  If a refined ball \(B'\subseteq B\) is also known to
contain a label, that label must be the unique label already in \(B\), and no
other label can occur by the same argument.
\end{proof}

This lemma is the correct interface for exact branch selection, root labeling,
spectral block counts, and discrete minimizers.  Without \(g\), the theorem
exports the cloud or multiset instead.

\subsection{Where relational output is essential}

The framework recommends relational output in the following cases.

\begin{itemize}
\item \textbf{Polynomial roots:} retain a multiset of isolating boxes and
multiplicities; label a simple root only with an isolation gap.
\item \textbf{Symmetric spectra:} ordered eigenvalues are stable, but bases
inside repeated eigenspaces are not; export blocks or projectors with a gap.
\item \textbf{Analytic branches:} retain all Newton/Puiseux candidates near a
singular point; choose one only with branch-separating data.
\item \textbf{Minimization:} the infimum value is stable without uniqueness;
a coherent point needs quantitative uniqueness or strong convexity.
\item \textbf{Covers and descent:} retain all finite local lifts and eliminate
dead ends by viability instead of selecting one globally in advance.
\end{itemize}

Relationality is not merely defensive.  It often yields a more invariant
statement and makes independent implementations comparable.
